\section{Finite completion by positive interpolation}
\label{sec:certificates}
The remaining degrees require a bound on the whole square, including a
strict comparison arbitrarily close to each maximizing corner. Both finite
methods use the same principle: interpolate a smooth signed sum with positive
weights, bound its interpolation error, and only then maximize over signs.
For any finite real vector $(h_\ell)$ the identity
\[
 \sum_\ell|h_\ell|=\max_{\varepsilon_\ell\in\{-1,1\}}
                         \sum_\ell\varepsilon_\ell h_\ell
\]
provides the signed sums. Each is smooth whenever the individual
kernels or cardinals are smooth. The signs realizing the maximum may
change inside a leaf; a bound valid for every complete signed sum
continues to hold across all such changes.

\begin{lemma}[Positive interpolation bounds]\label{cert:positive}
Let $F(z)=\max_{\sigma}f_\sigma(z)$, where the maximum is finite.
On a rectangle $R=[a_0,a_1]\times[b_0,b_1]$, suppose
$|(f_\sigma)_{aa}|\leq M\omega_a^2$ and
$|(f_\sigma)_{bb}|\leq M\omega_b^2$ for every $\sigma$.
If $W_z\geq F(z)$ at its vertices, then throughout $R$,
\begin{equation}\label{cert:peano}
 F(a,b)\leq I_RW(a,b)+\frac M2\{
 \omega_a^2(a-a_0)(a_1-a)+\omega_b^2(b-b_0)(b_1-b)\}.
\end{equation}
Here $I_R$ is bilinear interpolation. For polynomial cardinals with tensor
Bernstein coefficients $c_{\ell,ij}$ on $R$, the analogous bound is
\begin{equation}\label{cert:bernstein}
 \sum_\ell|L_\ell(z)|\leq
 \sum_{i,j}\left(\sum_\ell|c_{\ell,ij}|\right)B_i^n(s)B_j^n(t),
 \qquad (s,t)\in[0,1]^2.
\end{equation}
\end{lemma}
\begin{proof}
The one-dimensional linear interpolation remainder is an integral against
a Peano kernel of one sign, whose mass is $(a-a_0)(a_1-a)/2$.
Let $I_a,I_b$ denote the two linear interpolation operators. The identity
\[
 f_\sigma-I_aI_bf_\sigma
   =(f_\sigma-I_af_\sigma)+I_a(f_\sigma-I_bf_\sigma)
\]
separates the two errors. The first is bounded by
$M\omega_a^2(a-a_0)(a_1-a)/2$; positivity of $I_a$ bounds the second
by $M\omega_b^2(b-b_0)(b_1-b)/2$. At the four vertices,
$f_\sigma\leq F\leq W$, so positivity also gives
$I_aI_bf_\sigma\leq I_RW$. This proves \eqref{cert:peano} for
every signed sum, with an error independent of $\sigma$, and taking
their maximum proves the assertion.
Equation~\eqref{cert:bernstein} follows from the same positivity and
partition of unity of the Bernstein basis, after the triangle inequality.
\end{proof}

\subsection{One polynomial certificate for the small degrees}
For $2\leq n\leq40$ the direct certificate uses the corrected cardinal
kernel of Lemma~\ref{lem:cardinal-kernel}. Outward rational enclosures of
its coefficients are converted to tensor Bernstein form on the square.
Binary midpoint subdivision uses exact de Casteljau averages. On each leaf
all coefficient absolute sums in \eqref{cert:bernstein} must be strictly
less than a validated lower endpoint $L_c\leq\Lambda_c(n)$.
A leaf containing a designated corner exempts only its vertex coefficient:
that coefficient is exactly the polynomial value there, with absolute sum
$\Lambda_c(n)$; every other coefficient sum must be less than $L_c$.
Write $C=\Lambda_c(n)$, let $w_*$ be the exceptional
vertex weight, and let $D_R$ be the largest validated coefficient sum
among the other indices. The corner predicate gives $D_R<L_c\leq C$,
so the upper bound on this leaf is
\[
 Cw_*+D_R(1-w_*)=C-(C-D_R)(1-w_*).
\]
Here $w_*$ is a product of endpoint Bernstein weights, and $w_*<1$
at every point other than the designated vertex. This proves the
required strict comparison however close the point is to that vertex.
At the vertex itself the cardinal formula supplies the exact value $C$.
On an ordinary leaf the largest validated coefficient sum is already
strictly less than $L_c$, giving the comparison everywhere on the leaf.

The partition is specified by a rooted binary tree: every internal node
splits its closed rectangle at the midpoint of one coordinate, and both
children must be present. Starting with $[-1,1]^2$, this reconstructs a
cover of the entire square. Shared edges belong to both adjacent leaves;
the same inequalities hold there. The designated corner vertices are
exactly the two nonnode corners in Theorem~\ref{thm:main-location}.
Their node-preserving reflection gives their common value $C$.
Thus the leaf predicates prove both the maximum and its exact equality
set on the full square.

All coefficient enclosures are reconstructed by the checker. The identities
\[
 T_{r+1}(x)=2xT_r(x)-T_{r-1}(x),\qquad
 s^r=\sum_{i=r}^n\frac{\binom{i}{r}}{\binom nr}B_i^n(s)
\]
reduce construction and basis conversion to exact integer and rational
operations once the cosine arguments are bracketed. Alternating arctangent
series in Machin's identity bracket $\pi$; alternating cosine series and
the Chebyshev recurrence bracket the grid values. Integer operations round
outwards at every interval step. The checker reconstructs the whole
partition and its coefficients. Successful certificates are finite witnesses.
The accompanying execution record identifies which degrees were newly rerun
through the common implementation.

The degree-one case is explicit. Its nodes
are $(1,1),(1,-1),(-1,0)$, with cardinals
$(1+x+2y)/4$, $(1+x-2y)/4$, and $(1-x)/2$. Hence
\[
 \lambda_1(x,y)=\max\{1,|y|+(1-x)/2\},
\]
whose maximum is two, attained only at $(-1,\pm1)$.

\subsection{The phase certificate}
For $41\leq n\leq319$ the common envelope is enclosed on
$\mathcal T=\{(a,b):a,b\geq0,\ a+b\leq1\}$, with phases
$(\alpha,\beta)=(\pi a,\pi b)$ and $k=n/(n+1)$.
In the notation of \eqref{ul:primitive}, put
\[
 f_\varepsilon(a,b)=\frac1{2\mathsf N}
       \sum_{\Gamma^-}\varepsilon_{rs}H_{rs}(\pi a,\pi b).
\]
Then $S_n(\pi a,\pi b)=\max_\varepsilon f_\varepsilon(a,b)$,
and $|f_\varepsilon|\leq S_n$. Every signed board sum has frequencies
at most $\pi,\pi k$ in these coordinates. A genuine global norm $M$
and vertex uppers therefore give
\begin{align}\label{cert:phase-bound}
 S_n(\pi a,\pi b)&\leq q_R(a,b),\\
 q_R&=I_RW+\frac{M\pi^2}{2}
 \{(a-a_0)(a_1-a)+k^2(b-b_0)(b_1-b)\},\notag\\
 \max_Rq_R&\leq V_R+\theta_RM,
 \quad V_R=\max_zW_z,\quad
 \theta_R=\frac{\pi^2}{8}\{(a_1-a_0)^2+k^2(b_1-b_0)^2\}.\notag
\end{align}
Trigonometric Bernstein inequalities supply
$|(f_\varepsilon)_{aa}|\leq\pi^2M$ and
$|(f_\varepsilon)_{bb}|\leq\pi^2k^2M$ from the global norm.
The final line of \eqref{cert:phase-bound} uses positivity of the
bilinear weights and the elementary bound
$(a-a_0)(a_1-a)\leq(a_1-a_0)^2/4$, and similarly for $b$.
Thus $V_R+\theta_RM$ is a simple certified upper for the whole leaf.
The complete midpoint tree covers the triangle; a rectangle
is discarded only if its lower-left coordinate sum exceeds one.
Tangencies and shared boundaries remain in the cover.

The norm premise is itself certified. Let $M_*$ be the global real-phase
supremum of $S_n$. The folding symmetries established in
Section~\ref{sec:argmax} reduce this supremum to $\mathcal T$.
Compactness supplies a maximizing point, and every signed sum has
absolute value at most $M_*$. Consequently the derivative argument is
valid initially with this true, possibly unknown, norm. If all leaves
are numerical and $\theta_R<1$, a leaf containing that point gives
$M_*\leq V_R+\theta_RM_*$. Absorbing its last term and then taking
the largest possible leaf ratio yields
\begin{equation}\label{cert:tree-norm}
 M_*\leq\max_R\frac{V_R}{1-\theta_R}.
\end{equation}
Alternatively the quarter-grid points
$E=\{(i/4,j/4):i,j\geq0,\ i+j\leq4\}$ and their upper maximum $V_E$
give
\begin{equation}\label{cert:grid-norm}
 M_*\leq\frac{V_E}{1-\pi^2/64}<\frac{784}{663}V_E.
\end{equation}
Indeed the ten quarter-squares cover $\mathcal T$; outside vertices fold
back to $E$ by $(a,b)\mapsto(1-a,1-b)$, a phase symmetry.
Each square has side $1/4$, so its Peano error is at most
$\pi^2(1+k^2)M_*/128\leq\pi^2M_*/64$. Hence
$M_*\leq V_E+\pi^2M_*/64$, proving \eqref{cert:grid-norm}; the
rational factor follows from $\pi<22/7$. Either construction gives a
validated $M\geq M_*$ to insert in every subsequent leaf test.
For \eqref{cert:tree-norm}, every leaf meeting the triangle must
supply its numerical vertex bounds. A tree using an analytic far leaf
uses the independent grid norm \eqref{cert:grid-norm}.

On a leaf the certified upper for $q_R$ must be below $L_c$ on its far
part. On its actual near part $\rho\leq1/2$, at a noncorner center,
the required inequality is $q_R-\delta_R<L_c$, with the appropriate
proved reflection floor $\delta_R$. The floor's radius is
$\min\{1/2,\pi(a_1+b_1)\}$; it is used only where its physical guard
holds. To make the guard explicit, write $p_-\leq\pi\leq p_+$ for the
positive rational bracket. Since $\rho=\pi(a+b)$ on $\mathcal T$,
closed rational proxies for the two parts are
\[
 R\cap\mathcal T\cap\{a+b\leq1/(2p_-)\},\qquad
 R\cap\mathcal T\cap\{a+b\geq1/(2p_+)\}.
\]
The first contains every actual near point, and the second every actual
far point. Their overlap retains the seam $\rho=1/2$. Requiring the
appropriate strict comparison on each nonempty proxy is sufficient;
the defect floor is invoked in the physical implication only at actual
near points. Its certified radius can be evaluated outwards as
$r_R=\min\{1/2,p_+(a_1+b_1)\}$ and substituted in
\eqref{ul:finite-floors}. Every noncorner odd center has an eligible
projection, so \eqref{ul:defect} gives there
\[
 \lambda_n=S_n-\Delta_{JK}\leq q_R-\delta_n(r_R)<L_c
                                                   \leq\Lambda_c(n).
\]
On the actual far part $\lambda_n\leq S_n\leq q_R<L_c$.
Actual corners are covered by a direct descent witness.

The local descent uses \eqref{ul:squared-peano}. Take
$U=\pi^2/32$, $V=9\pi^2/512$, and outward actual-function values
$W_x,W_y,W_{xy}$ at $(\sqrt{2U},0),(0,\sqrt{2V}),
(\sqrt{2U},\sqrt{2V})$. With $u'=U/\pi^2$, $v'=V/\pi^2$, set
\[
 q=\min\left\{\frac{L_c-W_x}{u'},\frac{L_c-W_y}{v'},
                         \frac{L_c-W_{xy}}{u'+v'}\right\}.
\]
For the other three vertices the losses relative to the exact corner
value are at least $L_c-W_x$, $L_c-W_y$, and $L_c-W_{xy}$.
They are therefore at least $(q/\pi^2)U$, $(q/\pi^2)V$, and
$(q/\pi^2)(U+V)$, respectively. At the origin the value is exactly
$\Lambda_c(n)$. Apply \eqref{ul:squared-peano} with $c=q/\pi^2$
to obtain
\[
 \begin{gathered}
 \lambda_n(-\cos(A/n),-\cos(B/(n+1)))
 \leq\Lambda_c(n)-\frac{g}{2\pi^2}(A^2+B^2),\\
 g=q-\frac{M\pi^4}{6}\max\{u',k^4v'\}.
 \end{gathered}
\]
Thus the strict test $g>0$ proves positive quadratic loss off the
origin. The rectangle is $0\leq A\leq\pi/4$,
$0\leq B\leq3\pi/16$, which contains the inward half diamond
$A+B\leq1/2$. Complete signed cardinal sums are even in each inward
variable and have frequencies at most $1,k$; their absolute values are
bounded by the same global $M$, since $\lambda_n\leq S_n$ in every
chart. These are precisely the premises of the squared-variable
interpolation argument. The parity reflection transfers the witness to
the other actual corner.

From degree256 the analytic estimate \eqref{ul:corner256} supplies
the same strict descent on the required diamond. An analytic far leaf
is permitted from degree192 only when its entire triangle intersection
satisfies $\rho\geq\pi/2$. Equation~\eqref{ul:far192} then gives
$S_n<\Lambda_c(n)-3/100$ throughout that closed intersection.

To complete the implication, choose an odd chart for any point of the
physical square and fold its absolute phases into $\mathcal T$.
The tree gives a leaf containing those phases. Its far comparison,
noncorner near comparison, or actual-corner descent excludes the point
strictly, except when it is the origin of an actual corner chart.
Those origins are the two nonnode corners already identified by parity.
Their exact common value and the closed chart cover prove the asserted
maximum and equality set, including every chart tie and subdivision
boundary.

\subsection{Evidence and arithmetic}
The retained proof inventory covers every integer $2\leq n\leq319$;
Section~\ref{sec:argmax} covers the rest. For the phase range all280
primary/fresh pairs, including the control degree511, were checked by
exact rational reconstruction of560 essential proofs. Native outward
endpoints come from the frozen Arb/FLINT/MPFR/GMP implementation;
finite binary64 endpoints are converted losslessly to rationals for the
cover predicates. The arithmetic implementation, compiler and native
interfaces are trusted components of this computer-assisted proof.

The primary and fresh passes ran separately on the production machine.
Locally all essential proof/account/source links were checked. Complete
histories for two randomly selected degrees,199 and238, were inspected,
and fresh native verification was run at both degrees with distinct
review. These sample replays supplement
the checks of all essential proofs and links. Full unsampled
histories remain on the production machine under the selective-retention
policy. Source pins, witness identities, checker commands and the exact
new small-degree run inventory are in the companion reproducibility notes.
They permit regeneration; a durable public release remains to be arranged.
This evidence boundary is part of the theorem's computational scope.
